\documentclass[10pt,a4paper]{amsart}
\usepackage{amsmath,amssymb,amsthm,graphicx}
\usepackage{verbatim}
\newenvironment{reference}{\comment}{\endcomment}
%\newenvironment{reference}{}{}
\newenvironment{slogan}{\comment}{\endcomment}
\newenvironment{history}{\comment}{\endcomment}

% For commutative diagrams we use Xy-pic
\usepackage[all]{xy}

% We use 2cell for 2-commutative diagrams.
\xyoption{2cell}
% Theorem environments.
%
\theoremstyle{plain}
\newtheorem{theorem}[subsection]{Theorem}
\newtheorem{proposition}[subsection]{Proposition}
\newtheorem{lemma}[subsection]{Lemma}

\theoremstyle{definition}
\newtheorem{definition}[subsection]{Definition}
\newtheorem{example}[subsection]{Example}
\newtheorem{exercise}[subsection]{Exercise}
\newtheorem{situation}[subsection]{Situation}

\theoremstyle{remark}
\newtheorem{remark}[subsection]{Remark}
\newtheorem{remarks}[subsection]{Remarks}

\numberwithin{equation}{subsection}

% Macros
%
\def\lim{\mathop{\mathrm{lim}}\nolimits}
\def\colim{\mathop{\mathrm{colim}}\nolimits}
\def\Spec{\mathop{\mathrm{Spec}}}
\def\Hom{\mathop{\mathrm{Hom}}\nolimits}
\def\Ext{\mathop{\mathrm{Ext}}\nolimits}
\def\SheafHom{\mathop{\mathcal{H}\!\mathit{om}}\nolimits}
\def\SheafExt{\mathop{\mathcal{E}\!\mathit{xt}}\nolimits}
\def\Sch{\mathit{Sch}}
\def\Mor{\mathop{\mathrm{Mor}}\nolimits}
\def\Ob{\mathop{\mathrm{Ob}}\nolimits}
\def\Sh{\mathop{\mathit{Sh}}\nolimits}
\def\NL{\mathop{N\!L}\nolimits}
\def\CH{\mathop{\mathrm{CH}}\nolimits}
\def\proetale{{pro\text{-}\acute{e}tale}}
\def\etale{{\acute{e}tale}}
\def\QCoh{\mathit{QCoh}}
\def\Ker{\mathop{\mathrm{Ker}}}
\def\Im{\mathop{\mathrm{Im}}}
\def\Coker{\mathop{\mathrm{Coker}}}
\def\Coim{\mathop{\mathrm{Coim}}}

% Boxtimes
%
\DeclareMathSymbol{\boxtimes}{\mathbin}{AMSa}{"02}

%
% Macros for moduli stacks/spaces
%
\def\QCohstack{\mathcal{QC}\!\mathit{oh}}
\def\Cohstack{\mathcal{C}\!\mathit{oh}}
\def\Spacesstack{\mathcal{S}\!\mathit{paces}}
\def\Quotfunctor{\mathrm{Quot}}
\def\Hilbfunctor{\mathrm{Hilb}}
\def\Curvesstack{\mathcal{C}\!\mathit{urves}}
\def\Polarizedstack{\mathcal{P}\!\mathit{olarized}}
\def\Complexesstack{\mathcal{C}\!\mathit{omplexes}}
% \Pic is the operator that assigns to X its picard group, usage \Pic(X)
% \Picardstack_{X/B} denotes the Picard stack of X over B
% \Picardfunctor_{X/B} denotes the Picard functor of X over B
\def\Pic{\mathop{\mathrm{Pic}}\nolimits}
\def\Picardstack{\mathcal{P}\!\mathit{ic}}
\def\Picardfunctor{\mathrm{Pic}}
\def\Deformationcategory{\mathcal{D}\!\mathit{ef}}
\usepackage{cleveref}
\crefname{originalreference}{original reference}{original references}
\Crefname{originalreference}{Original reference}{Original references}
\makeatletter
\newlabel{lemma-dominate-by-normalized-blowing-up}{{0BBT}{source}{}{sourcepointer-1}{}}
\newlabel{lemma-dominate-by-normalized-blowing-up@cref}{{[originalreference][0][]0BBT}{[1][0][]source}{}{}{}}
\makeatother
\crefname{originalreference}{original reference}{original references}
\Crefname{originalreference}{Original reference}{Original references}
\makeatletter
\newlabel{lemma-exact-sequence}{{0AXF}{source}{}{sourcepointer-2}{}}
\newlabel{lemma-exact-sequence@cref}{{[originalreference][0][]0AXF}{[1][0][]source}{}{}{}}
\makeatother
\crefname{originalreference}{original reference}{original references}
\Crefname{originalreference}{Original reference}{Original references}
\makeatletter
\newlabel{lemma-rational-propagates}{{0BG0}{source}{}{sourcepointer-3}{}}
\newlabel{lemma-rational-propagates@cref}{{[originalreference][0][]0BG0}{[1][0][]source}{}{}{}}
\makeatother
\crefname{originallemma}{original lemma}{original lemmas}
\Crefname{originallemma}{Original lemma}{Original lemmas}
\makeatletter
\newlabel{limits-lemma-modifications}{{0B3X}{source}{}{sourcepointer-4}{}}
\newlabel{limits-lemma-modifications@cref}{{[originallemma][0][]0B3X}{[1][0][]source}{}{}{}}
\makeatother
\crefname{originallemma}{original lemma}{original lemmas}
\Crefname{originallemma}{Original lemma}{Original lemmas}
\makeatletter
\newlabel{limits-lemma-modifications-properties}{{0BFN}{source}{}{sourcepointer-5}{}}
\newlabel{limits-lemma-modifications-properties@cref}{{[originallemma][0][]0BFN}{[1][0][]source}{}{}{}}
\makeatother
\title{A maximal modification makes all local rings rational, and normalized blowups preserve this}
\author{The Stacks Project Authors}
\date{stacks-project:resolve.tex:local-snapshot-20261012; GNU FDL 1.2 or later}
\begin{document}
\maketitle
\noindent\textit{Editorial note.} All delivered mathematical text is the author's own, copied from the pinned Stacks Project chapter file resolve.tex; no statement or proof step was rewritten, reordered or removed.\par\medskip
\noindent\textit{Editorial note.} The counted proof is the author's complete argument: the choice of a normal modification maximizing the length of the first cohomology and the comparison of cohomology along further modifications, the spreading of modifications of the local rings to the ambient scheme with the maximality argument proving rationality, and the domination by normalized blowups with the propagation of rationality to the points above.\par\medskip
\noindent\textit{Editorial note.} The label lemma-reduce-to-rational names the counted statement, and the proof environment immediately following it in the source is delivered verbatim as its proof.\par\medskip
\noindent\textit{Editorial note.} The counted statement is carried by the author's own lemma environment in the chapter source.\par\medskip
\noindent\textit{Editorial note.} Pre-open gate: python operations/agent-authoring/tools/precard\_gate.py --source resolve.tex --label lemma-reduce-to-rational --start 2086 --end 2136 --local-id 28091 returned RESULT=OK (L28091 free, occupancy FREE, public interval clear). Pre-assembly tag self-check: canonical tag resolved and not present in reference/public-items.jsonl.\par\medskip
\noindent\textit{Editorial note.} The statement and proof are copied from the pinned chapter file; the cited inputs are kept as external pointers with their canonical Stacks tags.\par\medskip
\noindent\textit{Allowed standard prerequisites (editorial).}\begin{itemize}
\item Exact sequence relating the cohomology of a modification to that of a further one. Same chapter, the exact sequence used to compare the first cohomology of two modifications and to express the vanishing of the first direct image. Supplies the identification of the first cohomology groups of a modification and of a further modification dominating it.
\item Spreading a modification of a local ring to a modification of the ambient scheme. Limits, the lemmas spreading a modification of a local ring along a localization to a modification of the ambient scheme and checking properness. Supplies the modification of the maximal scheme restricting to a given modification of the local ring at a closed point.
\item Domination by normalized blowups and propagation of rationality along local extensions. Same chapter, the lemma dominating a modification by normalized blowups at closed points and the lemma that rationality propagates along essentially finite local extensions. Supply the sequence of normalized blowups dominating the maximal modification and the preservation of rationality at the points above.
\end{itemize}\medskip
\ifdefined\hypertarget\hypertarget{sourcepointer-1}{}\fi
\noindent\textit{Source reference.} Original reference 0BBT: sources/stacks-resolve/pinned-original-source.tex, source lines 1277--1277. Prerequisite pointer invoked at proof line(s) 2131; it names a standard fact of the same chapter that the counted argument uses, and every step of the counted proof built on it is retained verbatim in the delivered body.\par\medskip
\ifdefined\hypertarget\hypertarget{sourcepointer-2}{}\fi
\noindent\textit{Source reference.} Original reference 0AXF: sources/stacks-resolve/pinned-original-source.tex, source lines 1901--1901. Prerequisite pointer invoked at proof line(s) 2103; it names a standard fact of the same chapter that the counted argument uses, and every step of the counted proof built on it is retained verbatim in the delivered body.\par\medskip
\ifdefined\hypertarget\hypertarget{sourcepointer-3}{}\fi
\noindent\textit{Source reference.} Original reference 0BG0: sources/stacks-resolve/pinned-original-source.tex, source lines 2035--2035. Prerequisite pointer invoked at proof line(s) 2133; it names a standard fact of the same chapter that the counted argument uses, and every step of the counted proof built on it is retained verbatim in the delivered body.\par\medskip
\ifdefined\hypertarget\hypertarget{sourcepointer-4}{}\fi
\noindent\textit{Source reference.} Original lemma 0B3X: sources/stacks-chow/tags.txt, source lines 13980--13980. Prerequisite pointer invoked at proof line(s) 2113; it names a standard fact of a sibling chapter that the counted argument uses, and every step of the counted proof built on it is retained verbatim in the delivered body.\par\medskip
\ifdefined\hypertarget\hypertarget{sourcepointer-5}{}\fi
\noindent\textit{Source reference.} Original lemma 0BFN: sources/stacks-chow/tags.txt, source lines 14391--14391. Prerequisite pointer invoked at proof line(s) 2117; it names a standard fact of a sibling chapter that the counted argument uses, and every step of the counted proof built on it is retained verbatim in the delivered body.\par\medskip
\begin{lemma}
\label{lemma-reduce-to-rational}
Let $(A, \mathfrak m, \kappa)$ be a local normal Nagata domain
of dimension $2$. If reduction to rational singularities is possible
for $A$, then there exists a finite sequence of normalized blowups
$$
X = X_n \to X_{n - 1} \to \ldots \to X_1 \to X_0 = \Spec(A)
$$
in closed points such that for any closed point $x \in X$
the local ring $\mathcal{O}_{X, x}$ defines a rational singularity.
In particular $X \to \Spec(A)$ is a modification and $X$
is a normal scheme projective over $A$.
\end{lemma}

\begin{proof}
We choose a modification $X \to \Spec(A)$ with $X$ normal
which maximizes the length of $H^1(X, \mathcal{O}_X)$.
By Lemma \ref{lemma-exact-sequence}
for any further modification $g : X' \to X$ with $X'$ normal
we have $R^1g_*\mathcal{O}_{X'} = 0$ and
$H^1(X, \mathcal{O}_X) = H^1(X', \mathcal{O}_{X'})$.

\medskip\noindent
Let $x \in X$ be a closed point. We will show that $\mathcal{O}_{X, x}$
defines a rational singularity. Let $Y \to \Spec(\mathcal{O}_{X, x})$
be a modification with $Y$ normal. We have to show that
$H^1(Y, \mathcal{O}_Y) = 0$. By
Limits, Lemma \ref{limits-lemma-modifications}
we can find a morphism of schemes $g : X' \to X$ which is an isomorphism
over $X \setminus \{x\}$ such that $X' \times_X \Spec(\mathcal{O}_{X, x})$
is isomorphic to $Y$. Then $g$ is a modification as it is proper by
Limits, Lemma \ref{limits-lemma-modifications-properties}.
The local ring of $X'$ at a point of $x'$ is either isomorphic
to the local ring of $X$ at $g(x')$ if $g(x') \not = x$ and
if $g(x') = x$, then the local ring of $X'$ at $x'$ is isomorphic
to the local ring of $Y$ at the corresponding point. Hence we see
that $X'$ is normal as both $X$ and $Y$ are normal. By maximality
we have $R^1g_*\mathcal{O}_{X'} = 0$ (see first paragraph). Clearly
this means that $H^1(Y, \mathcal{O}_Y) = 0$ as desired.

\medskip\noindent
The conclusion is that we've found one normal modification $X$
of $\Spec(A)$ such that the local rings of $X$ at closed points all define
rational singularities. Then we choose a sequence of normalized
blowups $X_n \to \ldots \to X_1 \to \Spec(A)$ such that $X_n$
dominates $X$, see Lemma \ref{lemma-dominate-by-normalized-blowing-up}.
For a closed point $x' \in X_n$ mapping to $x \in X$ we can apply
Lemma \ref{lemma-rational-propagates} to the ring map
$\mathcal{O}_{X, x} \to \mathcal{O}_{X_n, x'}$
to see that $\mathcal{O}_{X_n, x'}$ defines a rational singularity.
\end{proof}

\end{document}
